\chapter{Slots and Content Projection}

Components are most reusable when they can accept external structure. In Relax, \emph{slots} are the mechanism for content projection.

\paragraph{What problem are slots solving?}

Without slots, a component can only vary by \emph{data} (props/state). But many UI building-blocks need to vary by \emph{structure}.
Examples:

\begin{itemize}
  \item \textbf{A Card} needs a customizable header/body/footer.
  \item \textbf{A TableRow} needs a caller-provided cell.
  \item \textbf{A Layout} needs to accept "whatever children" and place them inside a container.
\end{itemize}

In other frameworks, you might see this as:
\href{https://developer.mozilla.org/en-US/docs/Web/API/Web_components/Using_templates_and_slots}{Web Components \texttt{<slot>}},
\href{https://react.dev/learn/passing-props-to-a-component#passing-jsx-as-children}{React children},
or \href{https://vuejs.org/guide/components/slots.html}{Vue slots}.

Relax takes a deliberately small approach: model the slot as a \emph{temporary marker VNode}, then erase it with a deterministic tree rewrite.
This gives us content projection without making mounting or patching "slot-aware".

\paragraph{Tiny contract (keep this in your head)}

\begin{itemize}
  \item \textbf{Input}: A component renders a VNode tree that may contain SLOT markers (created by \texttt{hSlot()}).
  \item \textbf{Rewrite}: If the component received external children, \texttt{fillSlots()} rewrites the first slot position(s) into fragments containing chosen content.
  \item \textbf{Output}: After rewriting, the tree contains only normal VNodes (elements, fragments, text, components). No slots remain.
\end{itemize}

Relax’s slot model is intentionally minimal:

\begin{itemize}
  \item there is exactly one slot ``kind'' (no named slots)
  \item there is no runtime slot registry
  \item slots are compiled away into normal VDOM shape before mount/patch does serious work
\end{itemize}

That constraint keeps the renderer small and makes the behavior easy to specify with tests.

\section{Diagram: component render \textrightarrow slot rewrite \textrightarrow normal VDOM}
Slots are easiest to understand as a short rewrite pass that happens \emph{between} component render and mounting/patching:

\begin{center}
\begin{tikzpicture}[
  node distance=9mm,
  box/.style={draw,rounded corners,fill=RelaxLight,inner sep=5pt,align=center},
  arrow/.style={-Latex,thick}
]
  \node[box] (render) {component \texttt{render()}\\returns VNode tree\\(may contain SLOT markers)};
  \node[box, right=14mm of render] (flag) {\texttt{didCreateSlot()}?\\(module flag)};
  \node[box, right=14mm of flag] (fill) {\texttt{fillSlots(vdom, externalChildren)}\\mutates tree in-place};
  \node[box, below=of fill] (out) {output tree contains\\no SLOT markers\\(slots rewritten to FRAGMENT)};

  \draw[arrow] (render) -- (flag);
  \draw[arrow] (flag) -- node[above]{yes} (fill);
  \draw[arrow] (fill) -- (out);
\end{tikzpicture}
\end{center}

After this rewrite, slots disappear as a concept: mount/patch just see fragments and children arrays.

Slots in the VDOM layer are intentionally small and test-driven:

\begin{itemize}
  \item VNode marker: \texttt{DOM\_TYPES.SLOT} built by \texttt{hSlot()} in \texttt{src/h.ts}
  \item Slot filling algorithm: \texttt{fillSlots(vdom, children)} in \texttt{src/slots.ts}
  \item Specs: \texttt{src/\_\_tests\_\_/slots.test.ts} and \texttt{src/\_\_tests\_\_/component-slots.test.ts}
\end{itemize}

\section{Slot VNodes: \texttt{hSlot} and the SLOT marker}

A slot is represented as a VNode whose \texttt{type} is \texttt{DOM\_TYPES.SLOT}.

In \texttt{src/h.ts}:

\begin{itemize}
  \item \texttt{DOM\_TYPES} includes \texttt{SLOT: 'slot'}.
  \item \texttt{hSlot(children = [])} returns \texttt{\{ type: SLOT, children \}}.
  \item \texttt{hSlot} also sets a module-level flag \texttt{hSlotCalled = true}.
\end{itemize}

That flag is read via \texttt{didCreateSlot()} and reset via \texttt{resetDidCreateSlot()}.

\subsection{Why track \texttt{didCreateSlot}?}

Relax uses this flag as a cheap "did the render tree contain any slots?" signal.
The component renderer (\texttt{src/component.ts}) uses it:

\begin{itemize}
  \item if a slot was created during \texttt{render()}, it calls \texttt{fillSlots(vdom, externalChildren)}
  \item then it resets the flag
\end{itemize}

This avoids doing a full traversal to \emph{search} for slots when there are none.

\paragraph{How the optimization is wired (\texttt{src/component.ts})}

The component render pipeline is:

\begin{enumerate}
  \item call user \texttt{render()} to produce a VNode tree
  \item if any \texttt{hSlot()} call happened while rendering (module flag), run \texttt{fillSlots(vdom, externalChildren)}
  \item reset the module flag so the next render starts clean
\end{enumerate}

This is literally the code:

\begin{verbatim}
if (didCreateSlot()) {
  fillSlots(vdom, this.#children)
  resetDidCreateSlot()
}
\end{verbatim}

\subsection{Authoring model: what the user writes vs. what the runtime sees}

It helps to keep two mental views:

\begin{itemize}
  \item \textbf{Author view}: you "render a slot" and "pass children".
  \item \textbf{Runtime view}: you render a placeholder VNode of type \texttt{SLOT}, then run a rewrite that replaces it.
\end{itemize}

Here's a simple author-level example:

\begin{lstlisting}[language=JavaScript]
// A component with a slot position
const Card = defineComponent({
  render() {
    return h('section', { class: 'card' }, [
      h('header', {}, [h('h2', {}, ['Title'])]),
      h('main', {}, [hSlot([h('em', {}, ['Default body'])])]),
    ])
  },
})

// Call site provides external children
const app = h(Card, {}, [h('p', {}, ['Hello from the outside'])])
mountDOM(app, document.body)

// After slot rewrite, patch/mount only sees normal children under <main>
// (slot marker is gone)
\end{lstlisting}

Even if you never think about slots again, this explains why they're cheap: the renderer doesn't need special mount logic.

So slot filling is a \emph{conditional rewrite pass} executed only when a slot marker was created.

The tests explicitly check this optimization:

\begin{itemize}
  \item \texttt{when the component has slots, the "fillSlot()" function is called}
  \item \texttt{when the component has no slots, the "fillSlot()" function is not called}
\end{itemize}

(See \texttt{src/\_\_tests\_\_/component-slots.test.ts}).

\section{fillSlots(): projecting external children into a component}

The slot filling algorithm lives in \texttt{src/slots.ts}:

\begin{verbatim}
fillSlots(vdom, children)
\end{verbatim}

At a high level:

\begin{itemize}
  \item It traverses a VDOM tree.
  \item It finds SLOT nodes.
  \item It replaces a SLOT node with a FRAGMENT node containing either:
    \begin{itemize}
      \item external content (children passed into the component), or
      \item default slot content (children attached to the SLOT marker).
    \end{itemize}
\end{itemize}

\subsection{Contract}

\begin{itemize}
  \item Input: a rendered VDOM subtree \texttt{vdom}, and an array of \texttt{children} coming from the component call-site.
  \item Output: it mutates the tree in-place.
  \item Post-condition: every SLOT marker encountered is rewritten into a FRAGMENT node.
\end{itemize}

\subsection{A more precise contract: invariants you can rely on}

If you're debugging a slots issue, hand-wavy descriptions aren't enough. The tests assume these invariants:

\begin{itemize}
  \item \textbf{No slot markers survive}: after the rewrite, \texttt{DOM\_TYPES.SLOT} must not appear in the tree that gets mounted/patched.
  \item \textbf{One component = one filling context}: a component only fills slots within the VDOM tree returned by its own \texttt{render()}.
  \item \textbf{External children are late-bound per render}: on every component render, the caller-provided children array is considered "the current external content".
  \item \textbf{Mount/patch remain slot-agnostic}: any observable behavior must reduce to normal fragment/child reconciliation.
\end{itemize}

This design is why slots are "easy" in Relax: they are implemented as a deterministic pre-processing pass.

\subsection*{Diagram: the traversal + rewrite policy}

\begin{center}
\begin{tikzpicture}[
  node distance=8mm,
  box/.style={draw,rounded corners,fill=RelaxLight,inner sep=5pt,align=left},
  arrow/.style={-Latex,thick}
]
  \node[box] (walk) {Walk the VNode tree\\(pre-order traversal)};
  \node[box, below=of walk] (slot) {If node is SLOT:\\choose content\\external vs default};
  \node[box, below=of slot] (rewrite) {Rewrite SLOT \textrightarrow FRAGMENT\\\texttt{children = chosenContent}};
  \node[box, below=of rewrite] (drop) {During parent's child-walk:\\drop empty fragments\\that came from empty slots};

  \draw[arrow] (walk) -- (slot);
  \draw[arrow] (slot) -- (rewrite);
  \draw[arrow] (rewrite) -- (drop);
\end{tikzpicture}
\end{center}

\paragraph{Why FRAGMENT?}

The renderer needs to splice \emph{zero, one, or many} children into a single position in the parent.
In Relax’s VDOM model, fragments are the standard encoding for ``a list of children with no wrapping element''.
By rewriting slots to fragments, everything downstream (mounting, patching, diffing) can treat the result as normal VDOM.

\subsection{Choosing between default and external content}

When \texttt{fillSlots} sees \texttt{vdom.type === SLOT}, it computes:

\begin{itemize}
  \item \texttt{defaultContent} from \texttt{vdom.children}
  \item \texttt{externalContent} from the passed \texttt{children}
\end{itemize}

Then it uses a simple rule:

\begin{quote}
If external content is provided, it wins. Otherwise the default content is used.
\end{quote}

In code, this is:

\begin{verbatim}
const content = externalContent.length > 0 ? externalContent : defaultContent
\end{verbatim}

This is exactly what \texttt{src/\_\_tests\_\_/slots.test.ts} asserts:

\begin{itemize}
  \item \texttt{set default content if no external content}
  \item \texttt{set the external content if provided}
\end{itemize}

\subsection{Slots with no content are removed}

If a slot has neither default nor external content, Relax removes it.

Implementation detail:

\begin{itemize}
  \item A SLOT node is converted into a FRAGMENT.
  \item If the chosen content is empty, \texttt{vdom.children = []}.
  \item During the parent traversal pass, empty fragments that were slots are dropped.
\end{itemize}

Test: \texttt{remove the slot if no content or default content}.

\paragraph{A subtle implementation detail}

	exttt{fillSlots()} itself doesn't have access to a parent pointer.
So it can’t ``remove the slot node'' directly.
Instead it uses a two-step approach:

\begin{enumerate}
  \item rewrite SLOT \textrightarrow empty FRAGMENT (\texttt{children = []})
  \item during the parent’s child-walk pass, drop empty fragments that were produced by empty slots
\end{enumerate}

That design keeps the algorithm simple and makes behavior local to the traversal.

\subsection{A subtle but crucial point: \texttt{children} is not "DOM children"}

In Relax, component call-sites pass \emph{VNode children}, not DOM nodes.
That means "slot content" participates in the same rules we've already learned:

\begin{itemize}
  \item strings/numbers passed as children are normalized into TEXT VNodes (by \texttt{h()} / child mapping)
  \item arrays are structurally meaningful (fragments flatten in some places; slot-default content has a special shape rule)
  \item children are diffed like any other VNode list after slot filling
\end{itemize}

If you come from Web Components: HTML \texttt{<slot>} projects "real DOM".
Relax is closer to "project a subtree description".
Conceptual reference: \href{https://developer.mozilla.org/en-US/docs/Web/API/Web_components/Using_templates_and_slots}{MDN \texttt{<slot>}}.

\subsection{A code-walk of \texttt{fillSlots} (in plain JavaScript)}

The real implementation is in \texttt{src/slots.ts}, but this version captures the exact intent and shape.
The key is: \textbf{slot rewrite is mutation}, not allocation of a second tree.

\begin{lstlisting}[language=JavaScript]
function fillSlots(vdom, externalChildren) {
  if (!vdom) return

  // 1) Slot node found: rewrite it into a fragment
  if (vdom.type === DOM_TYPES.SLOT) {
    const defaultContent = Array.isArray(vdom.children) ? vdom.children : []
    const passed = Array.isArray(externalChildren) ? externalChildren : []
    const content = passed.length > 0 ? passed : defaultContent

    vdom.type = DOM_TYPES.FRAGMENT

    // Empty slot => become empty fragment; parent traversal later drops it
    if (content.length === 0) {
      vdom.children = []
      return
    }

    // Shape rule: hSlot([defaultContent]) means `content` is [[...]]
    const first = content[0]
    if (Array.isArray(first) && content.length === 1) {
      vdom.children = content
    } else {
      vdom.children = [content]
    }

    return
  }

  // 2) Traverse children, but treat components as opaque boundaries
  if (Array.isArray(vdom.children) && vdom.type !== DOM_TYPES.COMPONENT) {
    const next = []
    for (const child of vdom.children) {
      fillSlots(child, externalChildren)

      // Optional cleanup: drop empty fragments produced by empty slots
      if (child && child.type === DOM_TYPES.FRAGMENT && (!child.children || child.children.length === 0)) {
        continue
      }
      next.push(child)
    }
    vdom.children = next
  }
}
\end{lstlisting}

If you've implemented compilers, this should feel familiar: it's a single-pass rewrite with a few invariants.

\paragraph{Important: the real implementation differs slightly from the teaching version}

In the actual \texttt{src/slots.ts}, external content is set as:

\begin{quote}
fragment children = \texttt{content}
\end{quote}

not \texttt{[content]}. This matters for patching.
If you accidentally wrap external content in another array, you'd introduce an extra nesting level, which would make the fragment children shape wrong.

Interpretation:
	extbf{default content has a special container shape; external content is already a list.}

\subsection{Shape detail: why \texttt{hSlot([defaultContent])} is special}

There’s a slightly odd but intentional shape rule documented in \texttt{src/slots.ts}:

\begin{quote}
\texttt{hSlot([defaultContent])} passes an array whose first element is itself an array.
\end{quote}

So default slot content is often authored like:

\begin{verbatim}
hSlot([ h('span', {}, ['Hello']) ])
\end{verbatim}

In this case the algorithm sees: ``content is an array with a single element, and that element is also an array''.
It preserves that container shape so that:\\
	extbf{a slot expands to exactly one fragment child in the parent}, matching how the tests assert the resulting tree:

\begin{verbatim}
expect(vdom.children).toEqual([hFragment(defaultContent)])
\end{verbatim}

\section{Why fillSlots skips component children during traversal}

The traversal logic in \texttt{fillSlots} explicitly ignores children whose type is \texttt{COMPONENT}.

This is important. Consider:

\begin{itemize}
  \item You are filling a slot for component \texttt{A}.
  \item Inside \texttt{A}'s subtree there might be a nested component \texttt{B} that also uses slots.
\end{itemize}

Relax wants \textbf{each component to fill its own slots} using its own external content.
If the parent traversal reached into component children, it could fill the wrong slots with the wrong content.

Another way to say this: \texttt{fillSlots} treats \texttt{COMPONENT} vnodes as \emph{opaque}. That matches the rest of the renderer: components are patched/mounted as black boxes, and only the component instance is responsible for producing/maintaining its internal subtree.

This is a major boundary rule:

\begin{quote}
Slot rewriting is done per-component and must not cross component vnode boundaries.
\end{quote}

\subsection{Why this boundary matters in real apps}

If the parent traversal rewrote nested component internals, you'd get extremely confusing coupling:

\begin{itemize}
  \item A parent could accidentally fill a child component's slot with the wrong content.
  \item A refactor that wraps content in an extra component would silently change slot behavior.
  \item Slot filling would become order-dependent across renders (hard to test, hard to reason about).
\end{itemize}

By treating components as tuples of \texttt{(props, children)} and letting each component own its own slot rewrite,
Relax preserves a clean mental model: \textbf{slots are part of the component API surface, not the global tree.}

It also aligns with how the renderer treats components elsewhere:
components are patched/mounted as black boxes whose internal VDOM is managed by the component instance.

\subsection{From-scratch example: why crossing the boundary would be a bug}

Imagine this:

\begin{lstlisting}[language=JavaScript]
const Inner = defineComponent({
  render() {
    return h('div', {}, [hSlot([h('em', {}, ['Inner default'])])])
  },
})

const Outer = defineComponent({
  render() {
    return h('section', {}, [
      // Outer has a slot too
      hSlot([h('span', {}, ['Outer default'])]),

      // And it renders Inner
      h(Inner, {}, []),
    ])
  },
})

// Call site passes ONE set of children to Outer
mountDOM(h(Outer, {}, [h('b', {}, ['Projected into Outer'])]), document.body)
\end{lstlisting}

If \texttt{fillSlots} descended into the \texttt{Inner} component vnode,
it could (incorrectly) fill Inner's slot with Outer's external children.
That would mean Inner's content depends on who rendered it (bad encapsulation).

Relax avoids this by making \texttt{COMPONENT} a hard boundary during traversal.

This behavior is tested in \texttt{src/\_\_tests\_\_/slots.test.ts}:

\begin{itemize}
  \item \texttt{ignores slots whose parent is a component (they are handled by the component)}
\end{itemize}

\section{Slots in real components (\texttt{component-slots.test.ts})}

The integration tests show what authoring looks like.

\subsection{Basic projection}

\texttt{can have a slot where external content can be inserted}:

\begin{itemize}
  \item Component renders \texttt{<div>\{slot\}</div>}.
  \item Call site provides \texttt{<span>Hello</span>}.
  \item The mounted DOM contains the span.
\end{itemize}

\subsection{Default content}

\texttt{when no external content is provided, the default content is used}.

This is the standard "slot fallback" semantics.

\subsection{Conditional slots}

\texttt{conditionally rendered slots} verifies that slots can appear/disappear depending on component state, and the patcher will reconcile accordingly.

This is where slots meet the patching story from Chapter 8: slot output is just fragments and normal children after filling.

\subsection{Updating slot content between renders}

\texttt{slot content updated between renders} has a container component that alternates between:

\begin{itemize}
  \item a single \texttt{<span>Hello</span>}
  \item or two children \texttt{<p>World</p><span>!</span>}
\end{itemize}

The important point: slot external content is treated like regular children arrays and diffed normally after filling.

Practically:

\begin{itemize}
  \item the container component re-renders and passes a different external children array
  \item the slotted component re-renders, sees a slot marker, and re-runs \texttt{fillSlots}
  \item after that rewrite, the patcher sees ordinary FRAGMENT/ELEMENT/TEXT children and reconciles them
\end{itemize}

\subsection{Nested slots inside lists}

The last integration scenario in \texttt{component-slots.test.ts} is important because it looks like ``normal app code'':
build a list, render repeated components, and project a different child into each one.

Mentally, it's just two loops:

\begin{itemize}
  \item The parent maps \texttt{items} to \texttt{TableRow} component vnodes.
  \item Each \texttt{TableRow} render tree contains a slot marker in a \texttt{<td>}.
\end{itemize}

In author-level code (simplified):

\begin{lstlisting}[language=JavaScript]
const TableRow = defineComponent({
  render() {
    return h('tr', {}, [
      h('td', {}, [hSlot([h('em', {}, ['(empty)'])])]),
    ])
  },
})

const Table = defineComponent({
  render() {
    const titles = ['One', 'Two']
    return h('table', {}, [
      h('tbody', {}, titles.map((t) =>
        // Each row receives different external children.
        h(TableRow, {}, [t])
      )),
    ])
  },
})
\end{lstlisting}

The slot rewrite story stays exactly the same:
each \texttt{TableRow} instance stores its own external children (here: \texttt{['One']} or \texttt{['Two']}) and fills the slot while rendering.

Senior rule of thumb:
\begin{quote}
If you can explain repeated slots-in-lists, you understand slots. Everything else is a smaller case.
\end{quote}

\section{Walking the integration tests in slow motion (event-loop + render pipeline)}

Slots can feel magical until you align them with the component pipeline from earlier chapters.
At a high level, mounting a component with slot children looks like this:

\begin{enumerate}
  \item \textbf{Mount begins} at the parent vnode.
  \item When the renderer hits a \texttt{COMPONENT} vnode:
    \begin{enumerate}
  \item create the component instance
  \item store external children array on the instance (conceptually: \texttt{this.\#children})
  \item call \texttt{render()} to get an internal subtree
  \item if \texttt{didCreateSlot()} was set by \texttt{hSlot()}, run \texttt{fillSlots(subtree, this.\#children)}
  \item mount the rewritten subtree to DOM
    \end{enumerate}
\end{enumerate}

In other words, slot filling is part of "turn render output into mountable VDOM".
The patcher never sees a slot marker.

This is why the spy tests in \texttt{component-slots.test.ts} are important: they confirm the optimization boundary is hooked up correctly.

\texttt{slot in a nested component} uses a table row component with a slot in a \texttt{<td>}.
The parent maps titles \texttt{['One','Two']} into multiple \texttt{TableRow} instances.

This shows that slot content can be:

\begin{itemize}
  \item a raw string child (which becomes a TEXT VNode via \texttt{h()} normalization)
  \item repeated many times across a list of nested components
\end{itemize}

\section{Takeaways}

\begin{itemize}
  \item Slots are recorded as a VNode marker produced by \texttt{hSlot()}.
  \item Slot filling is a tree rewrite: SLOT \(\rightarrow\) FRAGMENT with chosen content.
  \item External content wins over default content.
  \item Empty slots are removed by dropping empty fragments.
  \item Each component fills its own slots; parent traversal does not cross component boundaries.
\end{itemize}

\section{Walking the tests: what they prove}

Slots are specified in two layers: a pure tree-rewrite unit suite and a DOM-level integration suite.

\subsection{Unit semantics (\texttt{src/\_\_tests\_\_/slots.test.ts})}

\begin{description}
  \item[remove the slot if no content or default content] Proves the empty-slot removal behavior: a slot with no default and no external content is eliminated (via empty-fragment dropping).
  \item[set default content if no external content] Proves fallback behavior and the default-content shape rule. The assertion expects \texttt{[hFragment(defaultContent)]}.
  \item[set the external content if provided] Proves the precedence rule: external content overrides default content.
  \item[ignores slots whose parent is a component] Proves the traversal boundary rule: \texttt{fillSlots} must not rewrite nested component internals.
\end{description}

\subsection{Integration semantics (\texttt{src/\_\_tests\_\_/component-slots.test.ts})}

\begin{description}
  \item[can have a slot where external content can be inserted] Proves end-to-end projection into mounted DOM.
  \item[default content is used] Proves fallback behavior at DOM level.
  \item[slot is empty] Proves empty slots become empty output (no wrapper nodes appear).
  \item[conditionally rendered slots] Proves slots can appear/disappear and patch correctly when state changes.
  \item[slot content updated between renders] Proves that changing the external children array results in normal DOM reconciliation after slot rewrite.
  \item[slot in a nested component] Proves nested component usage with lists and primitive children (strings become text nodes via \texttt{h()} normalization).
  \item[fillSlots spy tests] Proves the \texttt{didCreateSlot()} optimization works: slot rewrite only runs when \texttt{hSlot()} was actually used.
\end{description}

\section{Online resources}

\begin{itemize}
  \item Web Components: \href{https://developer.mozilla.org/en-US/docs/Web/API/Web_components/Using_templates_and_slots}{MDN \texttt{<slot>} (conceptual comparison)}
  \item Vue: \href{https://vuejs.org/guide/components/slots.html}{Slots and content projection}
  \item React: \href{https://react.dev/learn/passing-props-to-a-component#passing-jsx-as-children}{Passing JSX as children}
  \item SolidJS (conceptual): \href{https://www.solidjs.com/docs/latest#children}{Children utilities}
\end{itemize}

\section{Common pitfalls (senior notes)}

\begin{itemize}
  \item \textbf{Mistaking slots for "render props"}: Relax slots are structural children, not callbacks. If you need code execution with parameters, pass functions as props.
  \item \textbf{Forgetting the rewrite step exists}: by the time \texttt{mountDOM/patchDOM} runs, there is no slot node. Debug by inspecting the post-rewrite VDOM.
  \item \textbf{Component-boundary confusion}: if your slot isn't filling, check where you placed \texttt{hSlot()} -- it must be inside the component that receives the children.
  \item \textbf{Default content shape surprises}: remember that \texttt{hSlot([defaultContent])} uses a container shape that becomes a single fragment child.
\end{itemize}

If you internalize one idea: \textbf{Relax slots are a simple, deterministic rewrite into fragments}. Everything else follows.

\section{Performance notes (why this is fast in practice)}

Two design choices keep this lightweight:

\begin{itemize}
  \item \textbf{Conditional traversal}: the \texttt{didCreateSlot()} flag lets the component skip traversal when there are no slots.
  \item \textbf{In-place rewrite}: \texttt{fillSlots} mutates existing vnode objects. No cloning of the full tree.
\end{itemize}

That matters because slot filling runs on every render where slots are present.
This is the same kind of micro-optimization you see in production UI frameworks: avoid whole-tree work unless a feature is actually used.

Reference (conceptual): \href{https://react.dev/learn/render-and-commit}{React render/commit mental model}.

\section{Exercises}

\begin{enumerate}
  \item Modify the \texttt{fillSlots} teaching implementation so it matches \texttt{src/slots.ts} exactly---especially the external-content shape.
  \item Write a small component pair \texttt{Outer/Inner} (like the boundary example) and predict what would go wrong if traversal crossed into components.
  \item Explain (in 3-5 sentences) why "slots compiled away into fragments" simplifies the patcher.
\end{enumerate}

\section{Chapter summary}

\begin{itemize}
  \item Relax implements slots as a VNode marker (\texttt{DOM\_TYPES.SLOT}) created by \texttt{hSlot()}.
  \item Slot behavior is a deterministic rewrite pass: SLOT is rewritten to FRAGMENT containing either external or default content.
  \item Component vnodes are traversal boundaries: each component owns slot filling for its own render output.
  \item After filling, mount and patch operate on normal VDOM only---no special slot logic downstream.
\end{itemize}
